\begin{proof}
\begin{enumerate}[label=\arabic*.]
\item Put $R=\mathbb Z_{(3)}$ and use the rational coordinates of
\eqref{maps:eq:F.33}.  Let $\mathcal P_R$ be the algebra of blocks
$\left(\begin{smallmatrix}a&b\\0&D\end{smallmatrix}\right)$,
where $a\in R$, $b\in R^{1\times5}$ and $D\in\operatorname{Mat}_5(R)$.
We first determine the actual integral algebra generated by $A,B$.
The following ordered words have $31$ coordinate columns in the order
$(0,j)$, $0\leq j\leq5$, followed by $(i,j)$,
$1\leq i,j\leq5$:
\[
\begin{gathered}
I,A,B,A^2,AB,BA,B^2,A^3,A^2B,ABA,AB^2,BA^2,\\
BAB,B^2A,B^3,A^3B,A^2BA,A^2B^2,ABA^2,AB^2A,\\
AB^3,BA^3,BA^2B,BAB^2,B^2A^2,B^3A,A^3BA,\\
A^3B^2,A^2BA^2,BA^3B,B^2A^3.
\end{gathered}
\]
Their coordinate determinant is $-32$.  Its adjugate divided by
$-32$ therefore expresses every elementary matrix of
$\mathcal P_R$ as an $R$-linear combination of these words.
Since $A,B\in\mathcal P_R$, this proves
$R[A,B]=\mathcal P_R$ with the specified coordinate map.

For a stable lattice $L_R$, the elementary idempotents split its
six coordinate fractional ideals; the matrices $E_{ij}$ for
$i,j\geq1$ make the last five equal.  Thus
\[
 L_R=3^aRe_0\oplus3^bR^5,\qquad b\geq a,
\]
where the inequality follows from the matrices $E_{0j}$.
The matrix $N=AB-I$ has square zero and nonzero Smith factors
$1,1,1$, while $A_5B_5-I$ has factors $1,1$.
In this basis of $L_R$, its upper-right row is multiplied by
$3^{b-a}$.  For $b>a$ its reduction modulo three has rank two,
contradicting the primitive rank-three cusp.  Hence $b=a$.
The actual scalar map $3^{-a}I$ identifies $L_R$ with $R^6$
and intertwines every exterior power of the two generators.
The given cusp operator has the same elementary divisors:
$BA-I=A^{-1}(AB-I)A$, and
$(BA)^{-1}-I=-(BA)^{-1}(BA-I)$.
This proves the asserted local marking for every stable lattice.

All arithmetic in this item is independently reproduced by
\path{F2_12_three_primary_fresh.py} in \path{code/continuation_V/};
its output includes the full word-coordinate matrix, its determinant,
and the exact Smith and invariant-basis certificates used below.
The argument from that determinant treats all lattices, without a
bound on their indices.

\item Write $v_i=g_i^{m_i}$ for the actual finite power vectors.
They are fixed by the homology matrices $M_i$.
The base-orbifold quotient in \cref{maps:src:F.2.7} shows that
simple connectivity forces $\gcd(m_1,m_2)=1$.
After the cusp meridian relation the two base generators are
inverse up to a fibre translation.  Their actions on the abelian
fibre lattice are therefore inverse.  Their linear orders five
and twelve force both actions to be the identity in this quotient.
The exact Smith factors of $[A^t-I\ B^t-I]$ are five units,
with first row zero, so its fibre coinvariant module over $R$
is $R$ through the primitive character $\ell=x_0$.

Set $a_i=\ell(v_i)$.  An arbitrary cusp clutch contributes a scalar
$\delta$ to $g_1+g_2=\delta z$.  The remaining relation matrix is
\[
 \begin{pmatrix}m_1&0&-a_1\\0&m_2&-a_2\\1&1&-\delta\end{pmatrix}.
\]
Its determinant is a unit for a simply connected total space.
Since $3\mid m_2$ and $3\nmid m_1$, reduction modulo three makes
$a_2$ a unit in $R$.  This conclusion includes arbitrary clutch
translations and twists.

\item We determine the actual second finite-cap restriction over $R$.
Choose a primitive integral invariant character $\alpha$ whose
$R$-span is $\ell$.  The stable sublattice
\[
 \Lambda'=\mathbb Z v_2\oplus\ker\alpha
\]
has index $|\alpha(v_2)|$, prime to three.  It gives an actual
finite torus cover; the affine action lifts because it preserves
$v_2$, preserves $\ker\alpha$, and translates by $v_2/m_2$.
The corresponding reduced cap $F'_2$ covers $F_2$ with the same
degree.  In its basis beginning with $v_2$, the action is
\[
 (u,y)\longmapsto(u+1/m_2,Dy).
\]
Thus $F'_2$ is the mapping torus of $D:T^5\to T^5$, and its
regular six-torus is the $m_2$-fold cover in the circle direction.
The coefficient change has determinant a unit in $R$; under the
scalar marking of item~1, the vertical cohomology action $D^t$
is the displayed $B_5$.

The matrices $D^t-I$ and $\bigwedge^2D^t-I$ have no nonunit
three-primary Smith factors.  Their Wang sequences therefore make
$H^q(F'_2;R)$ free for $q\leq3$.  Pullback to the generic torus
is injective in these degrees, because transfer kills its kernel
by $m_2$ and the source is free.  Deck transformations of
$F'_2\to F_2$ act by translations on the generic torus cover.
Their action there is trivial on cohomology, so the injectivity
just proved makes their action on $H^q(F'_2;R)$ trivial as well.
Transfer for this prime-to-three cover now identifies
$H^q(F_2;R)$ with $H^q(F'_2;R)$, compatibly with their actual
restriction maps.  In particular the degree-two and degree-three
restrictions have zero kernel.

Put $T=D^t$.  Its polynomial is $(x-1)\Phi_{12}(x)$, and
\[
 E=\Phi_{12}(T)=T^4-T^2+I,\qquad E^2=E,\qquad
 \sum_{j=0}^{11}T^j=12E.
\]
The image of this integral rank-one idempotent is the primitive
invariant line.  Hence the norm for every stated $m_2$ is $m_2E$.
The Wang sequence in degree two has quotient
$\operatorname{Inv}(\bigwedge^2T)$ and kernel
$\operatorname{coker}(T-I)$.  Restriction to the $m_2$-fold cover
projects identically onto the former group.  On the latter it is
the horizontal wedge with the norm: a horizontal cell lifts to
$m_2$ successive transported prisms.
Consequently the restriction image $J_2$ projects onto the vertical
invariants, and its intersection with the horizontal invariant
summand is $e_0\wedge m_2\operatorname{Inv}(T)$.
The actual horizontal inclusion induces the isomorphism
\[
 \frac{\operatorname{Inv}(T)}{m_2\operatorname{Inv}(T)}
 \ \xrightarrow{\ \sim\ }\\
 \frac{\operatorname{Inv}(\bigwedge^2V_R)}{J_2}
 \ \cong\ R/3^rR.
\]
No choice of lifts of the vertical Wang classes changes this
quotient or its kernel.  At the first finite cap, $m_1$ is a unit
in $R$, so transfer and averaging give the full invariant image
and zero kernel.  At the cusp, \cref{night:P1.11} gives injective
restrictions onto the invariant lattices in degrees one, two
and three, by dualizing its actual saturated-kernel collapse.

\item The integral specialization-sheaf approach is used in
\cite[Lemma 7.25 and Propositions 7.26--7.27, pp.~61--63]{AC26}
and \cite[section ``The integral homology'']{E26}.
Here are the coefficient maps and calculations for the present lattice.  Put $L_q=\bigwedge^qV_R$, $A_q=\bigwedge^qA$,
$B_q=\bigwedge^qB$, and let $j$ include the thrice-punctured base.
For the full invariant-extension sheaf $j_*L_q$ the oriented
cellular complex is
\[
\begin{gathered}
 J_{1,q}\oplus J_{2,q}\oplus J_{0,q}
 \xrightarrow{d_0}L_q^2\xrightarrow{\Delta_q}L_q,\\
 J_{1,q}=\ker(A_q-I),\quad J_{2,q}=\ker(B_q-I),\quad
 J_{0,q}=\ker(B_qA_q-I),\\
 d_0(j_1,j_2,j_0)=(j_1-j_0,j_2-A_qj_0),\\
 \Delta_q(a,b)=(A_q-I)a+(B_q-I)b.
\end{gathered}
\]
These are the actual base attachment maps: the two cusp paths
carry $j_0$ and $A_qj_0$, whose combined boundary is
$(B_qA_q-I)j_0=0$.  Primitive invariant bases and a unimodular
kernel basis of $\Delta_q$ give the following certificates:
\[
\begin{array}{c|c|c|c|c}
q&\operatorname{rank}H^0&\operatorname{rank}H^1&
 \operatorname{SNF}(d_0\text{ in }\ker\Delta_q)&
 \operatorname{SNF}(\Delta_q)\\\hline
1&1&0&1^6&1^5,2\\
2&0&0&1^{15}&1^{15}\\
3&0&3&1^{16},2&1^{19},2.
\end{array}
\]
Here $1^s$ denotes $s$ entries equal to one.  The row-coordinate
matrices and their unimodular HNF transformations are specified
by the preceding formulas and retained in the independent output.
In particular $H^0(j_*L_2)=H^1(j_*L_2)=0$,
$H^2(j_*L_1)=0$ over $R$, and $H^0(j_*L_3)=0$.

Write $\mathcal R_q=R^qf_*R$.  The actual stalk/fibre
identification is proper base change
\cite[Tag 09V5, Lemma 20.18.1]{StacksProject}.
The local maps of item~3 give
\[
 0\longrightarrow\mathcal R_2\longrightarrow j_*L_2
 \longrightarrow Q\longrightarrow0,
 \qquad Q_{p_2}=R/3^rR,
\]
with no other support.  Hence its connecting homomorphism gives
$H^1(\mathcal R_2)=R/3^rR$.
The degree-one quotient is also point supported, so
$H^2(\mathcal R_1)=H^2(j_*L_1)=0$.
The degree-three sheaf map is injective, so
$H^0(\mathcal R_3)=0$.

\item The only total-degree-three term in the Leray sequence
\cite[Tag 01F2, Lemma 20.13.4]{StacksProject} is therefore
$E_2^{1,2}=R/3^rR$.  It has no possible differential because
the base has cohomological dimension two.  Thus
$H^3(M;R)=R/3^rR$.
Universal coefficients
\cite[Theorem 3.2, p.~195]{Hatcher2002} identify this torsion
module with $\operatorname{Ext}_R(H_2(M;R),R)$: the quotient
$\operatorname{Hom}_R(H_3(M;R),R)$ is free and must vanish.
The elementary divisor decomposition over the PID $R$ proves
\eqref{maps:eq:F.34}.

\item For the standard lattice take
\[
 v_1=(3,0,0,0,0,0)^t,\qquad
 v_2=(-7,7,-7,7,0,7)^t,\qquad (m_1,m_2)=(5,12).
\]
They are fixed by $A^t,B^t$ respectively.
Their $\ell$-values $3,-7$ are units modulo the respective
multiplicities; applying $\ell$ to a fixed-point equation proves
that every nonidentity affine power is free.
The primitive cusp in item~1 has the integral marking of
\cref{maps:src:F.2.1}.  Glue the two finite quotient caps and
the index-one Kuhn cap to the marked torus bundle over a pair
of pants, with zero additional clutch translation.  The finite
boundary identification is the prescribed logarithmic fibre
translation; the cusp is attached in that integral marking.
These actual local models give a closed smooth oriented
manifold.  Its remaining van Kampen matrix is
\[
 \begin{pmatrix}5&0&-3\\0&12&7\\1&1&0\end{pmatrix},
\]
with determinant of absolute value one.  It is simply connected.

Over $\mathbb Q$ the finite images are full invariants.
Item~4 gives vanishing $H^1(j_*V_{\mathbb Q})$,
$H^0(j_*\bigwedge^2V_{\mathbb Q})$,
$H^1(j_*\bigwedge^2V_{\mathbb Q})$ and
$H^2(j_*V_{\mathbb Q})$, and the degree-three invariant group
is zero.  The five-term Leray sequence and $H^1(M;\mathbb Q)=0$
make the map from the one-dimensional $H^0(\mathcal R_1)$ to
$H^2(\mathbb P^1;\mathbb Q)$ an isomorphism.
Thus $b_2=b_3=0$.
Euler additivity gives $\chi(M)=6$ from the Kuhn cusp;
all other fibres have Euler number zero.
Poincar\'e duality therefore gives
$(1,0,0,0,4,0,0,0,1)$, as claimed.
\end{enumerate}
\end{proof}
